\section{Splink and the causal cost of uncertain entity identity}
\label{sec:exec24-splink}
A bank linking servicing records to later outcomes must distinguish a plausible record pair from a trustworthy longitudinal observation. Splink supplies a concrete candidate for probabilistic linkage, comparison diagnostics and graph-based entity construction. These operations can support a causal workflow, but a matching score does not identify a treatment effect, certify a customer identity or supply the joint uncertainty of an entire linkage map. The proposed use therefore preserves candidate pairs, their selection mechanism and the information clock before considering any hard entity identifier.~\citep{exec24_splink_sources,exec24_splink_contract}

\subsection{Selected artifact and the boundary of the evidence}
The inspected repository revision is dated 6 October 2026 at 09:43:41 UTC, before the client cutoff. Its project configuration declares Splink 5.0.0, Python at least 3.10 and below 4, and direct dependencies on DuckDB, PyArrow and SQLGlot. Separate optional groups cover plotting, graph operations and additional database engines. The selected license file contains MIT terms. PyPI metadata independently identifies a 5.0.0 distribution uploaded on 28 September 2026; equality of the version string does not establish that the distribution has every byte of the later inspected source revision. No package was installed or executed. A deployment decision still needs an exact distribution, resolved dependencies, notices, supported database profile and maintenance owner.~\citep{exec24_splink_sources}

The source review authenticates selected files against the complete cutoff tree. It follows prediction, parameter training, blocking, threshold clustering, constrained clustering and evaluation delegates. Historical README and package metadata remain discovery evidence rather than substitutes for these technical paths. The review does not qualify every comparison function, database dialect, cache, persistence mechanism, optional extension or advertised scaling claim. Original execution behind a narrow adapter is a candidate; a wholesale port would require independently maintaining comparison semantics, SQL generation and graph algorithms without a demonstrated benefit.~\citep{exec24_splink_sources,exec24_splink_contract}

\subsection{A posterior score is conditional on a model and population}
Write $M$ for the event that two records represent the same entity, $\gamma$ for their comparison pattern and $\lambda=P(M)$ for the prior over a declared pair population. Under the selected comparison model, and suppressing term-frequency adjustments, posterior odds take the form
\[
 O(M\mid\gamma)=\frac{\lambda}{1-\lambda}\prod_j\frac{m_j(\gamma_j)}{u_j(\gamma_j)},
 \qquad P(M\mid\gamma)=\frac{O(M\mid\gamma)}{1+O(M\mid\gamma)}.
\]
Here $m_j$ and $u_j$ describe comparison outcomes among matches and nonmatches. The factorization needs its conditional-dependence assumptions; it is not a consequence of storing fields in separate columns. The inspected prediction path adds base-two log weights to the prior log odds, converts the result to a probability with separate positive and negative branches for numerical stability, and filters by an inclusive weight threshold. The null comparison level has zero weight. Selected term-frequency adjustments alter the weight using observed frequencies and configured floors. These numerical conventions must survive conversion to another interface.~\citep{exec24_splink_sources}

An authored exact calculation shows why the prior population matters. A likelihood ratio of 100 gives posterior probability $100/199$, about 0.503, at prior 0.01, but $100/1099$, about 0.091, at prior 0.001. The same observed agreement can therefore cross a decision threshold when the pair population or prior changes. If two fields are exact duplicates of one underlying signal with likelihood ratio 10, the correct one-signal calculation at prior 0.01 is $10/109$. Treating the duplicated fields as conditionally independent instead gives $100/199$. More agreeing columns have not created more information. These are specified probability examples, not empirical calibration measurements for Splink.~\citep{exec24_splink_contract}

The inspected prior-estimation method counts pairs admitted by deterministic rules, divides by a supplied recall estimate and then divides by the total possible comparisons. Its documentation requests rules with very few false positives, and its code checks the recall range and some count consistency conditions. Recall remains an external assumption. With 40 admitted pairs among 10,000 possible pairs, supplied recall one half produces prior 0.008; supplied recall four fifths produces 0.005. The calculation cannot validate the recall or the rules' truth from the admitted count alone. Boundary warnings should not be promoted into evidence that an extreme prior is scientifically appropriate.~\citep{exec24_splink_sources,exec24_splink_contract}

\subsection{Training design, stopping rules and unobserved comparisons}
The pinned training guide recommends directly estimating the prior and nonmatch probabilities before training selected match probabilities. Its rationale is explicitly conditional: random pairs should be overwhelmingly nonmatching, while blocking enriches the match population for training. The source API samples records and evaluates eligible pairs in chunks; its optional per-level count stopping rule is described as a heuristic rather than a precision guarantee. Seeds and database support have their own constraints. Reused records induce shared-pair dependence, so the number of processed pairs alone is not an independent-sample size for a confidence statement. A public prototype must preserve sampled record identities, pair inclusion, seed, chunking, early stopping and effective comparison counts.~\citep{exec24_splink_sources,exec24_splink_contract}

Training blocking also changes which comparisons can be estimated. The selected session removes comparisons involving blocking columns and adjusts the session prior with factors associated with the blocking rule. The API defaults to fixed nonmatch probabilities, while the session prior can change; using a trained session prior to populate the main model is a separate option that defaults to false. Thus a general recommendation to estimate the population prior directly must not be misread as a claim that every internal EM prior is held fixed. Different training blocks can supply complementary estimates, but agreement among them does not prove independence, correct selection or population calibration.~\citep{exec24_splink_sources}

The inspected EM loop compares the largest absolute parameter change with a tolerance and also imposes an iteration limit. Its final logging statement uses convergence language after the loop, including the path that exhausts that limit. This is a static control-flow observation, not an observed package failure. A proposed adapter must report whether the criterion was actually met, retain the parameter history and distinguish iteration exhaustion from successful convergence. An authored sequence of maximum changes $0.10,0.05,0.04$ with tolerance 0.01 never meets the criterion, even though three iterations have completed. Small parameter change, when obtained, is still not proof of global optimum or posterior calibration.~\citep{exec24_splink_sources,exec24_splink_contract}

\subsection{Blocking recall and the causal observation mechanism}
Prediction uses selected blocking rules or a manually registered blocked-pair table. The latter has its own chunking restrictions. Chunked prediction processes the designated combinations and combines outputs; a DuckDB-specific path prunes source inputs to records present in selected blocked pairs. Such execution paths require matched semantic tests before cross-engine or chunking equivalence is claimed. A pair excluded before scoring cannot be recovered merely by lowering a later probability threshold. The evaluated-label path explicitly distinguishes positives missed by blocking from errors in the scoring model. That distinction should be mandatory in the product's result contract.~\citep{exec24_splink_sources}

Suppose a public fixture contains 100 true links, of which blocking admits 60. The scorer retains 54 true links and six false links. Its recall among admitted true links is 90 percent and its precision is 90 percent, but end-to-end recall is only 54 percent. A report that conditions on scored pairs hides 46 missed true links. Labels sampled only from accepted pairs cannot identify the omitted recall; the evaluation design must include a justified sample outside the block, a known public oracle or an explicitly bounded sensitivity analysis. Annotation effort and unresolved labels belong in the cost denominator.~\citep{exec24_splink_contract}

For a banking illustration, consider two treatment arms, each containing outcomes 100 and 0 USD. Their complete-data contrast is zero. If the linkage mechanism retains only the high outcome in the treated arm and only the low outcome in the control arm, the observed contrast becomes 100 USD despite perfect identity for the retained pairs. This is selection bias, not a false-merge error. The example is a finite construction and supplies no causal identification claim. It motivates separate diagnostics for treatment-dependent linkage availability, outcome-dependent blocking and use of post-treatment contact fields. A customer address updated after a servicing intervention cannot silently become baseline identity evidence.~\citep{exec24_splink_contract}

\subsection{Pairwise probabilities do not define a joint entity map}
The selected threshold-clustering method constructs connected components from retained edges. It does not require every pair within a component to exceed the threshold. With scores 0.96 for $A$--$B$ and $B$--$C$, but 0.10 for $A$--$C$, threshold 0.95 joins all three records in an authored graph calculation; all six edge orders give the same component. This deterministic behavior is useful for explaining transitivity, but it cannot make the three scores a coherent joint posterior over equivalence relations. If the first two numbers were probabilities under one common posterior, then transitivity and the union bound require
\[P(A\sim C)\geq P(A\sim B)+P(B\sim C)-1=0.92.\]
The score 0.10 violates that requirement. The example does not assert that a real dataset generates those scores; it shows why independent pair labels and a joint identity distribution are different objects.~\citep{exec24_splink_sources,exec24_splink_contract}

The reviewed alternative single-best-link routine constrains selected datasets to at most one member per cluster and exposes distinct tie policies. Such a source-uniqueness constraint can be valuable, but a local mutual-best rule is not automatically a maximum-total-weight assignment. In an authored two-by-two score matrix with rows $(0.90,0.80)$ and $(0.85,0.10)$, choosing the mutually highest pair first and then the remaining pair totals 1.00; the alternative cross assignment totals 1.65. This is an objective comparison, not an execution result or an assertion that total score is the correct banking loss. One-to-one assumptions themselves fail for legitimately repeated records unless their scope is explicitly declared.~\citep{exec24_splink_sources,exec24_splink_contract}

A causal consumer needs either a justified joint distribution over feasible linkage maps, an identified set with sensitivity bounds, or a refusal when identity information is inadequate. If a specified joint model assigns probabilities three quarters and one quarter to two maps yielding contrasts $20$ and $-20$ USD, the model-averaged contrast is 10 USD and the between-map variance is 300 USD squared. If each conditional analysis has variance 4 USD squared, the law of total variance gives 304 USD squared. Reporting only 4 discards the identity contribution. These weights are stipulated joint hypotheses; they cannot be manufactured by independently drawing the package's pairwise probabilities. Calibration and causal validity still require the joint model, observation law and downstream assumptions.~\citep{exec24_splink_contract}

\subsection{Reference disposition and a testable integration route}
The completed proposal choice is to retain Splink as a reference and conditional prototype candidate, without granting native integration. A deterministic exact-key baseline should establish the value of probabilistic matching. A pinned original worker is preferable to a wholesale port when it can preserve the selected database and comparison semantics. Alternative linkage or joint-identity models should be compared on the same known maps, missingness patterns, causal targets and full loss function. Better pair discrimination alone is not an adequate selection criterion.~\citep{exec24_splink_contract}

The proposed CLI/MCP request binds immutable input tables, source/record identifiers, allowed feature clocks, link type, candidate-generation rules, parameter and term-frequency snapshots, declared seeds, thresholds, tie behavior, engine identity and a resource budget. The response separates candidate pairs, excluded/unscored status, estimated probabilities, clustering rule, proposed entity maps, convergence diagnostics and uncertainty limitations. Failed, refused and unresolved requests stay visible. A result must not label a blocked-out pair as a calibrated nonmatch or an estimated component as certified identity. Arbitrary SQL and database credentials require their own scoped execution boundary; this source review has contacted no live bank or MCP endpoint.~\citep{exec24_splink_contract}

Before admission, a public/generated study should freeze planted entity maps, duplicate fields, correlated missingness, treatment-dependent linkage, conflicting source identifiers, bridge edges, tie permutations, blocked-out positives, iteration exhaustion and engine/chunking profiles. It should compare pair accuracy, entity consistency, causal estimand error, refusal and complete cost. The scientific owner controls the target and linkage law; the data owner controls clocks and source truth; the adapter owner preserves conversion and failure semantics; the platform and rights owners qualify the selected distribution and environment. Authentic bank validation follows a useful prototype and remains a separate obligation.~\citep{exec24_splink_contract}

An authored cost example stipulates 40 USD computation, 160 USD annotation, 200 USD engineering and 80 USD storage/review, totaling 480 USD. Six attempts for four logical requests, three of which produce accepted artifacts, cost 80 USD per attempt, 120 USD per request and 160 USD per accepted artifact. These are explanatory inputs, not prices or a Splink benchmark. Revisit the reference decision only after measured need, permissible artifact identities and predeclared semantic, uncertainty, failure and cost criteria support a bounded integration. A faster matching engine cannot supply missing identity truth or a universally calibrated causal analysis.~\citep{exec24_splink_contract}

The chapter's reuse choices therefore follow the information each component preserves. A catalogue locates and describes evidence; preparation and feature tools construct the declared population and history; validation records which conditions were actually checked; linkage retains uncertainty about the unit; and lineage or scheduling records how the work proceeded. An existing engine behind a narrow reviewed interface is the conditional preference where it serves that role. Explicit local files and finite oracles remain the small-case comparators, while native qualification determines whether the additional machinery earns its complete cost. No single catalogue status, successful job or matching score can substitute for the combined observation and causal argument.
